\chapter{Further Integrable Models}
\label{chap:further-integrable-models}

The previous two chapters solved the XXZ chain by two complementary methods.
The coordinate Bethe ansatz emphasized factorized scattering of magnons, while
the algebraic Bethe ansatz emphasized the Yang--Baxter algebra of the six-vertex
transfer matrix.  This chapter explains how the same ideas extend beyond the
nearest-neighbor XXZ chain.  The goal is not to give a second complete solution
of every model.  Instead, the purpose is to identify the structural data that
make each model solvable and to locate it within the mother-model organization
of these notes.

We discuss three classes of examples:
\begin{enumerate}[label=(\roman*)]
    \item the isotropic XXX chain, obtained as the rational limit of the XXZ
    chain;
    \item the one-dimensional Hubbard model, whose exact solution requires a
    nested Bethe ansatz because particles carry both charge and spin;
    \item long-range integrable models, represented by the Haldane--Shastry spin
    chain and the Calogero--Sutherland model.
\end{enumerate}
The recurring theme is that integrability is not synonymous with free
fermions.  In these models the many-body spectrum is controlled by nontrivial
quantization conditions, hidden symmetries, or special long-range interactions.
The excitations are still strongly constrained, but they need not be free
Bogoliubov quasiparticles.

\section{The XXX limit}
\label{sec:further-xxx-limit}

The XXZ Hamiltonian used in \cref{chap:coordinate-bethe-ansatz-xxz,chap:algebraic-bethe-ansatz}
was
\begin{equation}
\label{eq:further-xxz-hamiltonian-recall}
    H_{\rm XXZ}^{(+)}
    =
    J\sum_{j=1}^{L}
    \left(
        \sigma_j^x\sigma_{j+1}^x
        +\sigma_j^y\sigma_{j+1}^y
        +\Delta(\sigma_j^z\sigma_{j+1}^z-1)
    \right),
\end{equation}
with periodic boundary conditions.  The isotropic XXX chain is the special point
\begin{equation}
\label{eq:further-xxx-limit}
    \Delta\to1,
    \qquad
    H_{\rm XXX}^{(+)}
    =
    J\sum_{j=1}^{L}
    \left(
        \bm{\sigma}_j\cdot\bm{\sigma}_{j+1}-1
    \right).
\end{equation}
Equivalently, in terms of spin operators $\bm{S}_j=\bm{\sigma}_j/2$,
\begin{equation}
\label{eq:further-xxx-spin-form}
    H_{\rm XXX}^{(+)}
    =
    4J\sum_{j=1}^{L}
    \left(
        \bm{S}_j\cdot\bm{S}_{j+1}-\frac14
    \right).
\end{equation}
The constant is chosen so that the fully polarized reference state has zero
energy, as in the XXZ Bethe ansatz chapters.

\begin{remark}[Sign convention]
With $J>0$ in \cref{eq:further-xxx-limit}, a one-magnon state has
$\varepsilon(k)=4J(\cos k-1)\leq0$ relative to the ferromagnetic reference
state.  Many condensed-matter texts instead write the antiferromagnetic XXX
Hamiltonian as
$J_{\rm AF}\sum_j\bm{S}_j\cdot\bm{S}_{j+1}$ with $J_{\rm AF}>0$.
The two conventions differ by an overall sign and an additive constant.  The
Bethe equations are convention-independent; only the energy formula changes by
this normalization.
\end{remark}

\subsection{Rational R-matrix}
\label{subsec:further-rational-r-matrix}

The XXX chain is also obtained from the six-vertex $R$-matrix by taking the
rational limit.  In the hyperbolic convention of the algebraic Bethe ansatz
chapter,
\begin{equation}
    R_{\rm XXZ}(u)
    =
    \begin{pmatrix}
        \sinh(u+\eta)&0&0&0\\
        0&\sinh u&\sinh\eta&0\\
        0&\sinh\eta&\sinh u&0\\
        0&0&0&\sinh(u+\eta)
    \end{pmatrix},
    \qquad
    \Delta=\cosh\eta .
\end{equation}
Sending $u\mapsto\eta u$ and then taking $\eta\to0$ gives, up to an irrelevant
scalar normalization,
\begin{equation}
\label{eq:further-rational-R-matrix}
    R_{\rm XXX}(u)
    =u\,\id+P,
\end{equation}
where $P$ permutes the two spin-$1/2$ spaces:
\begin{equation}
\label{eq:further-permutation-operator}
    P\ket{\alpha}\otimes\ket{\beta}
    =
    \ket{\beta}\otimes\ket{\alpha}.
\end{equation}
This rational $R$-matrix satisfies the Yang--Baxter equation
\begin{equation}
\label{eq:further-rational-ybe}
    R_{12}(u-v)R_{13}(u-w)R_{23}(v-w)
    =
    R_{23}(v-w)R_{13}(u-w)R_{12}(u-v).
\end{equation}
It is regular because
\begin{equation}
\label{eq:further-rational-regularity}
    R_{\rm XXX}(0)=P.
\end{equation}
Therefore the logarithmic derivative of the rational transfer matrix gives a
nearest-neighbor Hamiltonian.  Using
\begin{equation}
    P_{j,j+1}=\frac12\left(\id+\bm{\sigma}_j\cdot\bm{\sigma}_{j+1}\right),
\end{equation}
one obtains the isotropic exchange interaction.

\begin{principle}[Trigonometric versus rational integrability]
The XXZ chain is controlled by the trigonometric six-vertex $R$-matrix.  The
XXX chain is controlled by its rational degeneration.  The solution method is
unchanged: monodromy matrix, transfer matrix, RTT algebra, Bethe vectors, Bethe
equations.  What changes is the analytic form of the scattering phase.
\end{principle}

\subsection{Bethe equations in rapidity form}
\label{subsec:further-xxx-bethe-equations}

The rational rapidity variable $\lambda$ is related to the one-magnon momentum
by
\begin{equation}
\label{eq:further-xxx-momentum-rapidity}
    \ee^{\ii k}
    =
    \frac{\lambda+\frac{\ii}{2}}{\lambda-\frac{\ii}{2}}.
\end{equation}
For $M$ magnons, the XXX Bethe equations are
\begin{equation}
\label{eq:further-xxx-bethe-equations}
    \left(
        \frac{\lambda_j+\frac{\ii}{2}}
             {\lambda_j-\frac{\ii}{2}}
    \right)^L
    =
    \prod_{\substack{\ell=1\\ \ell\neq j}}^{M}
    \frac{\lambda_j-\lambda_\ell+\ii}
         {\lambda_j-\lambda_\ell-\ii},
    \qquad
    j=1,\ldots,M .
\end{equation}
The corresponding two-body scattering amplitude is
\begin{equation}
\label{eq:further-xxx-two-body-S}
    S_{\rm XXX}(\lambda_j-\lambda_\ell)
    =
    \frac{\lambda_j-\lambda_\ell+\ii}
         {\lambda_j-\lambda_\ell-\ii}.
\end{equation}
Thus the many-body quantization condition still has the universal Bethe form
\begin{equation}
\label{eq:further-general-bethe-quantization}
    \ee^{\ii k_jL}
    =
    \prod_{\ell\neq j}S(k_j,k_\ell),
\end{equation}
but the scattering phase has become rational rather than trigonometric.

Using \cref{eq:further-xxx-momentum-rapidity}, one finds
\begin{equation}
\label{eq:further-xxx-cos-rapidity}
    \cos k
    =
    \frac{\lambda^2-\frac14}{\lambda^2+\frac14}.
\end{equation}
Therefore the energy in the normalization of \cref{eq:further-xxx-limit} is
\begin{equation}
\label{eq:further-xxx-energy}
    E
    =
    4J\sum_{j=1}^{M}(\cos k_j-1)
    =
    -2J\sum_{j=1}^{M}
    \frac{1}{\lambda_j^2+\frac14}.
\end{equation}
The total lattice momentum is
\begin{equation}
\label{eq:further-xxx-total-momentum}
    \ee^{\ii P}
    =
    \prod_{j=1}^{M}
    \frac{\lambda_j+\frac{\ii}{2}}
         {\lambda_j-\frac{\ii}{2}}.
\end{equation}

\subsection{Enhanced SU(2) symmetry}
\label{subsec:further-xxx-su2}

For generic $\Delta$, the XXZ chain conserves only the $U(1)$ charge
$S^z_{\rm tot}$.  At $\Delta=1$, the symmetry is enhanced to full spin
rotation symmetry:
\begin{equation}
\label{eq:further-xxx-su2-commutator}
    [H_{\rm XXX}^{(+)},S^\alpha_{\rm tot}]=0,
    \qquad
    S^\alpha_{\rm tot}=\frac12\sum_{j=1}^{L}\sigma_j^\alpha,
    \qquad
    \alpha=x,y,z .
\end{equation}
This enhancement has an important consequence.  Bethe states constructed above
the all-up pseudovacuum are usually highest-weight states of $SU(2)$ multiplets.
The remaining members of the multiplet are obtained by acting with
$S^-_{\rm tot}$.  This is why completeness questions in the XXX chain are
slightly more delicate than merely counting finite rapidity solutions of
\cref{eq:further-xxx-bethe-equations}.

\begin{remark}[Infinite rapidities]
In the XXX chain, acting with the global spin-lowering operator can be viewed as
adding a Bethe root at infinity.  This is a useful algebraic way to understand
how $SU(2)$ descendants fit into the Bethe-ansatz framework.
\end{remark}

\section{Hubbard model and nested Bethe ansatz}
\label{sec:further-hubbard-model}

The one-dimensional Hubbard model is the canonical example of an interacting
fermion model that is exactly solvable but not reducible to a single-species
magnon Bethe ansatz.  Its local degrees of freedom are
\begin{equation}
\label{eq:further-hubbard-local-space}
    \ket{0},
    \qquad
    \ket{\uparrow},
    \qquad
    \ket{\downarrow},
    \qquad
    \ket{\uparrow\downarrow},
\end{equation}
so each site carries both charge and spin.  The Hamiltonian on a periodic chain
is
\begin{equation}
\label{eq:further-hubbard-hamiltonian}
    H_{\rm Hub}
    =
    -t\sum_{j=1}^{L}\sum_{\sigma=\uparrow,\downarrow}
    \left(
        c_{j\sigma}^\dagger c_{j+1,\sigma}
        +c_{j+1,\sigma}^\dagger c_{j\sigma}
    \right)
    +U\sum_{j=1}^{L}n_{j\uparrow}n_{j\downarrow} .
\end{equation}
Here
\begin{equation}
    n_{j\sigma}=c_{j\sigma}^\dagger c_{j\sigma},
    \qquad
    c_{L+1,\sigma}=c_{1\sigma}.
\end{equation}
The hopping amplitude is $t$, and $U$ is the on-site interaction.

\subsection{Why a nested ansatz is needed}
\label{subsec:further-why-nested}

In the XXZ chain, the Bethe particles are down spins moving in an all-up
reference state.  There is only one species of quasiparticle, so every collision
is described by a scalar two-body phase.  In the Hubbard model, a particle
carries charge momentum and spin.  When two fermions scatter, their spin labels
can be rearranged.  The two-body scattering matrix therefore acts on an internal
spin space.  Diagonalizing the charge problem leaves a second spin problem to be
solved.  This second step is the nested Bethe ansatz.

The exact solution is parametrized by two sets of variables:
\begin{equation}
\label{eq:further-hubbard-variables}
    \{k_j\}_{j=1}^{N},
    \qquad
    \{\lambda_\alpha\}_{\alpha=1}^{M},
\end{equation}
where $N$ is the total number of electrons and $M$ is the number of down-spin
electrons.  The $k_j$ are charge momenta, while the $\lambda_\alpha$ are spin
rapidities.

\begin{principle}[Nested Bethe ansatz]
When the two-body scattering matrix carries an internal index, Bethe ansatz
becomes a two-level problem.  First one writes a coordinate Bethe ansatz for
charge motion.  Then one diagonalizes the internal spin scattering problem by a
second Bethe ansatz.  The final answer contains coupled Bethe equations for
charge and spin rapidities.
\end{principle}

\subsection{Lieb--Wu equations}
\label{subsec:further-lieb-wu-equations}

Introduce the dimensionless interaction parameter
\begin{equation}
\label{eq:further-hubbard-u-parameter}
    u=\frac{U}{4t}.
\end{equation}
For the Hamiltonian \cref{eq:further-hubbard-hamiltonian}, the Lieb--Wu Bethe
equations are
\begin{align}
\label{eq:further-lieb-wu-charge}
    \ee^{\ii k_j L}
    &=
    \prod_{\alpha=1}^{M}
    \frac{\sin k_j-\lambda_\alpha+\ii u}
         {\sin k_j-\lambda_\alpha-\ii u},
    \qquad
    j=1,\ldots,N,\\
\label{eq:further-lieb-wu-spin}
    \prod_{j=1}^{N}
    \frac{\lambda_\alpha-\sin k_j+\ii u}
         {\lambda_\alpha-\sin k_j-\ii u}
    &=
    \prod_{\substack{\beta=1\\ \beta\neq\alpha}}^{M}
    \frac{\lambda_\alpha-\lambda_\beta+2\ii u}
         {\lambda_\alpha-\lambda_\beta-2\ii u},
    \qquad
    \alpha=1,\ldots,M .
\end{align}
The energy and total momentum are
\begin{equation}
\label{eq:further-hubbard-energy-momentum}
    E=-2t\sum_{j=1}^{N}\cos k_j,
    \qquad
    P=\sum_{j=1}^{N}k_j
    \quad (\mathrm{mod}\ 2\pi),
\end{equation}
for the unshifted convention in \cref{eq:further-hubbard-hamiltonian}.  If one
uses the particle-hole symmetric Hamiltonian
\begin{equation}
\label{eq:further-shifted-hubbard}
    H'_{\rm Hub}
    =
    -t\sum_{j,\sigma}
    \left(
        c_{j\sigma}^\dagger c_{j+1,\sigma}
        +c_{j+1,\sigma}^\dagger c_{j\sigma}
    \right)
    +U\sum_j
    \left(n_{j\uparrow}-\frac12\right)
    \left(n_{j\downarrow}-\frac12\right),
\end{equation}
then the same Bethe roots give the shifted energy
\begin{equation}
\label{eq:further-shifted-hubbard-energy}
    E'=E-\frac{U}{2}N+\frac{U}{4}L .
\end{equation}

\begin{remark}[Why $\sin k$ appears]
The Hubbard equations are not obtained by simply replacing the XXZ rapidity by
$k$.  The charge dispersion is that of a tight-binding lattice fermion,
$-2t\cos k$, while the spin scattering depends on $\sin k$.  This mixture is a
characteristic feature of the Hubbard nested solution.
\end{remark}

\subsection{Half filling and the Mott gap}
\label{subsec:further-hubbard-half-filling}

At half filling, $N=L$.  For repulsive interaction $U>0$, the exact solution
shows that the charge sector is gapped while the spin sector remains gapless.
This is the one-dimensional Mott-insulating mechanism in its cleanest exact
form.  The Hilbert space still allows charge fluctuations, but the low-energy
sector separates into
\begin{equation}
\label{eq:further-hubbard-spin-charge-structure}
    \text{gapped charge excitations}
    \qquad\text{and}\qquad
    \text{gapless spin excitations}.
\end{equation}
For large positive $U/t$, virtual hopping produces the effective antiferromagnetic
Heisenberg chain
\begin{equation}
\label{eq:further-hubbard-large-U-heisenberg}
    H_{\rm eff}
    =
    J_{\rm ex}\sum_{j=1}^{L}
    \left(
        \bm{S}_j\cdot\bm{S}_{j+1}-\frac14
    \right),
    \qquad
    J_{\rm ex}=\frac{4t^2}{U},
\end{equation}
inside the subspace with one electron per site.  Thus the XXX chain appears as
the spin-sector low-energy limit of the repulsive half-filled Hubbard model.

\begin{example}[From Hubbard to Heisenberg]
Consider two neighboring sites with one electron on each site.  A hopping event
creates a virtual doubly occupied site and costs energy $U$.  A second hopping
event returns to the singly occupied subspace.  Second-order perturbation theory
therefore gives an exchange scale $t^2/U$.  Fermionic statistics select the
antiferromagnetic exchange, producing $J_{\rm ex}=4t^2/U$.
\end{example}

\subsection{Integrability structure}
\label{subsec:further-hubbard-integrability}

The Hubbard model is integrable, but its integrability is more subtle than that
of the six-vertex model.  The relevant $R$-matrix is not the elementary
six-vertex $R$-matrix; it is a higher structure compatible with two coupled
spin species.  Equivalently, one may view the model as having two intertwined
XX-type hopping problems connected by the on-site interaction.  The important
lesson for these notes is the following:
\boxedpar{%
The Hubbard model is solvable not because it is quadratic, but because its
many-body scattering is still factorized after accounting for the nested charge
and spin structure.}

\section{Long-range integrable models}
\label{sec:further-long-range-integrable-models}

The nearest-neighbor XXX, XXZ, and XYZ chains teach that local Hamiltonians can
be integrable when they arise from logarithmic derivatives of commuting transfer
matrices.  There is another important class of integrable models whose defining
feature is not nearest-neighbor locality, but a very special long-range
interaction.  The two standard examples are the Haldane--Shastry spin chain and
the Calogero--Sutherland particle model.

\subsection{Haldane--Shastry chain}
\label{subsec:further-haldane-shastry}

Place $L$ spin-$1/2$ degrees of freedom at equally spaced points on the unit
circle,
\begin{equation}
\label{eq:further-hs-roots-of-unity}
    z_j=\ee^{2\pi\ii j/L},
    \qquad
    j=1,\ldots,L .
\end{equation}
The Haldane--Shastry Hamiltonian is
\begin{equation}
\label{eq:further-haldane-shastry-hamiltonian}
    H_{\rm HS}
    =
    J\left(\frac{\pi}{L}\right)^2
    \sum_{1\leq j<\ell\leq L}
    \frac{\bm{S}_j\cdot\bm{S}_\ell-\frac14}
         {\sin^2\left(\frac{\pi(j-\ell)}{L}\right)} .
\end{equation}
The interaction decays as the inverse square of the chord distance.  Although
this Hamiltonian is long-ranged, it is more exactly solvable than a generic
nearest-neighbor interacting spin chain.

At half filling of down spins, an exact ground-state wave function can be written
in terms of the down-spin coordinates $z_1,\ldots,z_{L/2}$ as a lattice Jastrow
state,
\begin{equation}
\label{eq:further-hs-ground-state}
    \Psi_{\rm HS}(z_1,\ldots,z_{L/2})
    =
    \prod_{a<b}(z_a-z_b)^2
    \prod_{a=1}^{L/2} z_a .
\end{equation}
This is the spin-chain analogue of a Laughlin-type wave function.  The square
Jastrow factor enforces strong correlations between spin flips.

The elementary excitations are spinons.  They carry spin $1/2$ rather than spin
$1$, and they obey fractional exclusion statistics.  This is qualitatively
different from the magnons of the ferromagnetic Bethe ansatz, which are spin-flip
particles above a polarized vacuum.

\begin{remark}[Yangian symmetry]
The hidden symmetry of the Haldane--Shastry chain is a Yangian extension of
$SU(2)$.  This enlarged symmetry is responsible for large spectral degeneracies
and for the unusually simple spinon structure.  It is the long-range analogue of
the algebraic constraints imposed by the Yang--Baxter equation in nearest-neighbor
chains.
\end{remark}

\subsection{Calogero--Sutherland model}
\label{subsec:further-calogero-sutherland}

The Calogero--Sutherland model is a continuum many-body system on a circle with
inverse-square interactions.  A common normalization is
\begin{equation}
\label{eq:further-calogero-sutherland-hamiltonian}
    H_{\rm CS}
    =
    -\sum_{j=1}^{N}\frac{\partial^2}{\partial x_j^2}
    +
    g(g-1)\left(\frac{\pi}{L}\right)^2
    \sum_{1\leq j<\ell\leq N}
    \frac{1}
         {\sin^2\left(\frac{\pi(x_j-x_\ell)}{L}\right)} .
\end{equation}
The exact ground state has the Jastrow form
\begin{equation}
\label{eq:further-cs-ground-state}
    \Psi_0(x_1,\ldots,x_N)
    =
    \prod_{j<\ell}
    \left|
        \sin\frac{\pi(x_j-x_\ell)}{L}
    \right|^{g}.
\end{equation}
The excited states are obtained by multiplying this ground state by symmetric
polynomials.  More precisely, Jack polynomials provide the natural eigenbasis.
Thus the Calogero--Sutherland model is solvable by special-function and
representation-theoretic methods rather than by a nearest-neighbor coordinate
Bethe ansatz.

The coupling $g$ controls the statistical interaction between particles.  At
$g=0$ or $g=1$, the inverse-square potential vanishes, corresponding to free
bosonic or fermionic endpoints depending on the imposed symmetry.  For general
$g$, the model realizes fractional exclusion statistics in a continuum setting.

\subsection{Freezing trick and relation to spin chains}
\label{subsec:further-freezing-trick}

The Haldane--Shastry chain and the Calogero--Sutherland model are not unrelated.
In the strong-coupling limit of spin Calogero--Sutherland-type models, the
particle coordinates freeze near equally spaced classical equilibrium positions
on the circle.  The remaining low-energy dynamics is carried by the internal
spin variables.  This limiting procedure produces inverse-square spin-chain
Hamiltonians of Haldane--Shastry type.

Schematically,
\begin{equation}
\label{eq:further-freezing-trick-schematic}
    \text{spin Calogero--Sutherland particles}
    \xrightarrow[\text{coordinates freeze}]{\text{large coupling}}
    \text{Haldane--Shastry spin chain}.
\end{equation}
This relation is called the freezing trick.  It is another example of the
mother-model principle: a continuum integrable particle model contains a lattice
integrable spin chain as a controlled limit.

\section{Comparison of integrable structures}
\label{sec:further-comparison-integrable-structures}

The examples in this chapter illustrate several distinct meanings of exact
solvability.  They are summarized in \cref{tab:further-integrability-comparison}.

\begin{table}[htbp]
\centering
\small
\begin{tabularx}{\linewidth}{@{}L{0.22\linewidth}L{0.25\linewidth}Y@{}}
\toprule
Model & Solvability mechanism & Main spectral data \\
\midrule
XXX chain & Rational $R$-matrix; Bethe ansatz; enhanced $SU(2)$ symmetry & Rapidities satisfying rational Bethe equations \\
XXZ chain & Trigonometric six-vertex $R$-matrix; coordinate/algebraic Bethe ansatz & Rapidities satisfying trigonometric Bethe equations \\
XYZ chain & Elliptic eight-vertex $R$-matrix; Baxter $TQ$ relation & Elliptic spectral parameter and functional equations \\
Hubbard chain & Nested Bethe ansatz; factorized scattering with charge and spin variables & Charge momenta $k_j$ and spin rapidities $\lambda_\alpha$ satisfying Lieb--Wu equations \\
Haldane--Shastry chain & Inverse-square interaction; Yangian symmetry; spinon basis & Fractional-statistics spinons and exact Jastrow ground state \\
Calogero--Sutherland model & Inverse-square continuum interaction; Jack-polynomial eigenfunctions & Pseudomomenta and Jack-polynomial labels \\
\bottomrule
\end{tabularx}
\caption{Representative integrable models beyond the basic XXZ coordinate
Bethe ansatz.  The word ``integrable'' refers to different but related
structures: rational, trigonometric, or elliptic $R$-matrices; nested Bethe
ansatz; Yangian symmetry; or inverse-square special-function solvability.}
\label{tab:further-integrability-comparison}
\end{table}

The following hierarchy is useful:
\begin{equation}
\label{eq:further-integrability-hierarchy}
\begin{gathered}
    \text{free fermion solvability}
    \subset
    \text{Bethe/Yang--Baxter integrability}
    \subset
    \text{broader exact solvability},\\
    \text{but none of these inclusions should be interpreted as equality.}
\end{gathered}
\end{equation}
Free-fermion models are solvable because the Hamiltonian becomes quadratic.
Bethe-ansatz models are solvable because many-body scattering factorizes.
Long-range inverse-square models are solvable because the interaction is tuned
to special symmetry and polynomial structures.  Commuting-projector models,
which appear in the next part of the notes, are solvable for yet another reason:
the Hamiltonian decomposes into mutually compatible local constraints.

\section{What should be remembered}
\label{sec:further-what-to-remember}

The main points of this chapter are:
\begin{enumerate}[label=(\roman*)]
    \item The XXX chain is the rational limit of the XXZ chain.  Its Bethe
    equations are obtained by replacing trigonometric scattering data by rational
    scattering data.
    \item The XXX point has enhanced $SU(2)$ symmetry.  Bethe roots at infinity
    and highest-weight multiplets are natural consequences of this symmetry.
    \item The Hubbard model requires a nested Bethe ansatz because the particles
    carry both charge and spin.  Its solution is encoded in the coupled
    Lieb--Wu equations.
    \item At half filling and repulsive interaction, the Hubbard model has a
    gapped charge sector and a gapless spin sector.  In the large-$U$ limit, the
    low-energy spin dynamics reduces to an antiferromagnetic Heisenberg chain.
    \item Long-range inverse-square models such as the Haldane--Shastry and
    Calogero--Sutherland models are integrable for reasons different from the
    nearest-neighbor transfer-matrix construction.  Their exact wave functions
    and spectra are governed by Jastrow factors, Yangian symmetry, and special
    polynomials.
\end{enumerate}

\boxedpar{%
The common lesson is that exact solvability is a structural property.  Quadratic
fermions, Yang--Baxter algebras, nested scattering, inverse-square interactions,
and commuting projectors are different mechanisms, but each replaces the generic
many-body problem by a constrained algebraic or analytic problem.}
